\documentclass[11pt]{article}
\usepackage[a4paper,margin=25mm]{geometry}
\usepackage{fontspec}
\usepackage[hidelinks]{hyperref}
\setmainfont{Latin Modern Roman}
\pagestyle{plain}
\begin{document}
{\LARGE\bfseries ALIUS Bulletin - Guidelines for Interviews - 2022\par}
\vspace{1em}
\section*{Alius Bulletin}
\section*{*}
\section*{Guidelines For}
\section*{Interviews}
If you want to conduct an interview for ALIUS Bulletin, please read the following guidelines carefully.\par

Our aim for this Bulletin is to include at least 5 interviews.\par

New schedule system\par

For this Bulletin, we introduce a rolling publication system. Interviews are published whenever they are ready, i.e., when all stages of the publication process are fulfilled (see modus operandi and requirements below). All interviews are integrated at the end of the year in the annual Bulletin. Pages number and citation format are updated in the final publication.\par

Length of the interview\par

Each interview is expected to be of at least 3500 words of length (this includes the interviewer’s questions, the interviewee’s responses and the bibliography).\par

Content of the interview\par

You are not expected to carry out the interview in a journalistic style – asking short and non-technical questions. Think of this interview almost as a dialogue. If needed, do not hesitate to ask long and detailed questions. Bear in mind that interviewees do not necessarily have a lot of time available and might take many things for granted. Therefore, your job, as the interviewer, is to make sure that your questions summarize in an understandable manner key concepts and facts that will facilitate the reading for non-specialists. So, when deemed relevant, do not hesitate to write a few lines introducing key concepts and key data to contextualize a bit the questions you are asking. That being said, if you want to ask short questions, feel free to do so: write long questions only when the intelligibility of what you are saying requires it.\par

Generally speaking, the interviews are intended for scholars. Therefore, you are not expected to ask general audience questions as a science journalist would do. However, each interview should be easily understandable to any scholar. For example, an anthropologist reading an interview conducted with a neuroscientist should easily understand most of the discussion, and conversely, a neuroscientist reading an interview conducted with an anthropologist should easily understand most of the discussion.\par

We invite you to have a look at Fortier \& Friedman’s interview with Friston or Limanowski \& Millière’s interview with Metzinger in the second issue of the Bulletin, to have a concrete example of how questions can be phrased. However, these two models are just two illustrations of what can be done. Feel free to differ from these models: as long as you broadly follow the above recommendations, it will be perfect.\par

Modus operandi and requirements\par

1. Ask one of the editors of the Bulletin whether the researcher you want to interview fits well with the Bulletin’s editorial line.\par

2. Ask the researcher you want to interview whether he/she agrees in principle to be interviewed (when you do so, let him/her know about the general concept of the Bulletin, about the other researchers interviewed in previous issues of the Bulletin, specify that the interview is done in a written way, describe the modus operandi, the schedule, etc.).\par

3. You can now write up your questions. Feel free to discuss them with one of the editors if needed, but this is not required.\par

4. Next, send your questions to the interviewee, and when you do so, set up a deadline specifying when the interviewee is expected to send back his/her responses.\par

(5). (This step is optional) Once you receive the interviewee’s responses, and if time permits, you may be willing to send a second round of questions. That is, some of the interviewee’s responses may call for further questions. It’s up to you to decide to send the interviewee’s a new series of questions.\par

6. Once you have all the responses from the interviewee, you will next have to:\par

-Give a title to the interview (it is better to discuss this directly with the interviewee as you will need his/her final agreement for the title and the content of the interview).\par

-Write a short abstract (<200 words) and up to 5 keywords for referencing the interview (you will need the final agreement of the interviewee for these).\par

-If the interviewee discusses books or articles in the interview but does not provide the detailed reference, it is your job to insert a bibliography at the end of the interview where all the books and articles mentioned by you (in your questions) or the interviewee (in his responses) are referenced. For the citation style that you are expected to use, see step 7 below.\par

-Select 2 to 6 excerpts (depending on the length of the interview) highlighting some key points made by the interviewee and nicely epitomizing his/her ideas. These passages must be no longer than one or two sentences (to have a clearer idea as to how to choose these passages, see highlighted excerpts in the first issue of the Bulletin). Insert the 2, 3, 4 or 5, 6 selected excerpts at the very end of the text.\par

7. When the file is ready – when it includes the interview (questions and responses), a title, a bibliography and a few selected excerpts, send the whole file to the editors. Note that the text in the file must comply with the following requirements:\par

-Use Times New Roman, size 12.\par

-All your questions must appear in red; all the interviewee’s responses must appear in black.\par

-Citation style for the bibliography is the 6th edition of APA (if you use Zotero, this citation style can be downloaded at the following link: https://www.zotero.org/styles?q=id\%3Aapa).\par

-The Bulletin is written in American English and not in British English. If the text that you or the interviewee wrote is in British English, it will have to be turned into American English (e.g.: programme => program; realised => realized; etc.). This can be easily be done using Word Autocorrect. If yourself or the interviewee uses British lexicon or phrasings, that’s fine: you can keep it as it is. But typographic details must be changed to fulfill the rules of American English. In addition to the “s” transformed into “z”, keep in mind that in American English, abbreviations such as “e.g.” and “i.e.” are followed by a comma (for example: “i.e., the mind is…” and not “i.e. the mind is…”).\par

-Use the following quotation marks: “” (do not use «» or ‘’). If there is a quotation within a quotation, use the following quotation marks: ‘’. For example: “I know that the meaning he gives to ‘consciousness’ differs from…”\par

-Make sure that there is no space before punctuation marks (; : , . ! ?). For example it should be: …because of two reasons: first, it doesn’t… and not: …because of two reasons : first, it doesn’t…\par

-Dashes used for punctuation must be “em dashes” and not “en dashes” and they must not be separated by spaces (see: https://en.wikipedia.org/wiki/Dash). For example, it should be: …the concept of consciousness—thus defined—proves difficult… and not: …the concept of consciousness – thus defined – proves difficult…
(We cannot really expect interviewees to comply with all these requirements. So each interviewer will be in charge of making sure that the interviewee’s responses do comply with these requirements.)\par

-Before sending the file to the editors, take time to carefully go through the interview to make sure there is no typo.\par

8. Once they receive the text of the interview, the editors will forward it to a native English-speaking proofreader who will take care of checking grammar and linguistic correctness. Note that the job of proofreaders is not to correct typos. As indicated above, typos must be corrected before sending the file to the editors.\par

9. The editors will next send you the feedback provided by the proofreaders. You will be expected to integrate to the interview the changes suggested by the proofreaders.\par

10. Once all changes have been integrated, read the whole interview again very carefully to make sure that there is no typo left.\par

11. When the final version is ready, send it to the interviewee to get his/her final agreement for publication.\par

12. Once you have the interviewee’s agreement for publication, send the final version of the interview to one of the editors. The interview will then be ready to be published!\par

(13.) If the interviewee doesn’t agree to publish the interview as it is and wants some changes to be made, make these changes and then send the revised version of the piece to the editors; they will forward it once more to the proofreaders to make sure there are no grammatical/linguistic problems. Once proofreaders have gone through the interview, we will proceed as described in step 9 above.\par

For any questions, feel free to ask editors (here below)\par

Editorial team\par

Daniel Friedman (danielarifriedman@gmail.com); Matthieu Koroma (matthieu.koroma@gmail.com) ; George Fejer (gyorgyf@live.com)\par

\end{document}
